\documentclass[../main.tex]{subfiles}
\graphicspath{{\subfix{../images/}}}
\begin{document}

Im morphologischen Kasten (Abbildung \ref{fig:morphkasten}) werden die Teilprobleme der Mindmap als Parameter in den Zeilen aufgeführt. In den Spalten stehen die zugehörigen Lösungsprinzipien. Durch die Kombination je einer Lösung pro Zeile entstehen Gesamtlösungen. Aus der Vielzahl möglicher Kombinationen wurden drei Gesamtlösungen ausgewählt, die in der Abbildung farblich hervorgehoben sind. Dabei wurden nur Kombinationen betrachtet, die die Festforderungen aus Kapitel 2 grundsätzlich erfüllen können. Gesamtlösung 3 (Aluminium) verletzt die Festforderung der Funkdurchlässigkeit und wird daher in Kapitel 7 verworfen. Die Gesamtlösungen 1 und 2 werden anschließend in der Nutzwertanalyse verglichen.

\usetikzlibrary{calc}

% ---------- Farben ----------
\definecolor{mkHead}{RGB}{0,51,102}
\definecolor{mkParam}{RGB}{214,224,238}
\definecolor{mkCell}{RGB}{242,245,249}
\definecolor{mkText}{RGB}{30,30,30}
\definecolor{mkA}{RGB}{0,113,178}
\definecolor{mkB}{RGB}{213,92,0}
\definecolor{mkC}{RGB}{0,158,115}

% ---------- Raster ----------
\newcommand{\mkN}{7}      % Anzahl der Spalten (ohne Parameterspalte)
\newcommand{\mkM}{8}      % Anzahl der Parameterzeilen
\newcommand{\mkP}{3.4}    % Breite der Parameterspalte in cm
\newcommand{\mkH}{1.5}    % Zeilenhöhe in cm

\newcommand{\mkSetup}{%
  \pgfmathsetmacro{\mkW}{(\linewidth/1cm-\mkP)/\mkN}%
  \pgfmathsetmacro{\mkTW}{\mkW-0.1}%
  \pgfmathsetmacro{\mkPTW}{\mkP-0.2}%
}

% ---------- Zellstile ----------
\tikzset{
  mkbase/.style={
    rectangle,
    minimum height=\mkH cm,
    inner sep=0.8pt,
    draw=white,
    line width=1.2pt,
    font=\sffamily\tiny,
    text=mkText
  },
  mkkopf/.style={mkbase, minimum width=\mkW cm, text width=\mkTW cm,
    fill=mkHead, text=white, align=center, font=\sffamily\scriptsize\bfseries},
  mkkopfp/.style={mkbase, minimum width=\mkP cm, text width=\mkPTW cm,
    fill=mkHead, text=white, align=center, font=\sffamily\scriptsize\bfseries},
  mkparam/.style={mkbase, minimum width=\mkP cm, text width=\mkPTW cm,
    fill=mkParam, font=\sffamily\scriptsize\bfseries,
    execute at begin node=\raggedright},
  mkzelle/.style={mkbase, minimum width=\mkW cm, text width=\mkTW cm,
    fill=mkCell, execute at begin node=\centering}
}

% ---------- Makros ----------
\newcommand{\mkkopfp}[2]{%
  \node[mkkopfp] (z#1x0) at (0,{-#1*\mkH}) {#2};}
\newcommand{\mkparam}[2]{%
  \node[mkparam] (z#1x0) at (0,{-#1*\mkH}) {#2};}
\newcommand{\mkkopf}[3]{%
  \pgfmathsetmacro{\mkXX}{\mkP/2+(#2-0.5)*\mkW}%
  \node[mkkopf] (z#1x#2) at (\mkXX,{-#1*\mkH}) {#3};}
\newcommand{\mkzelle}[3]{%
  \pgfmathsetmacro{\mkXX}{\mkP/2+(#2-0.5)*\mkW}%
  \node[mkzelle] (z#1x#2) at (\mkXX,{-#1*\mkH}) {#3};}

% \mkLoesung{Farbe}{Linienstil}{rel. x-Versatz}{Spaltenliste}
\newcommand{\mkLoesung}[4]{%
  \pgfmathsetmacro{\mkDX}{#3*\mkW}%
  \foreach \sp [count=\zl] in {#4}{%
    \coordinate (mkpt\zl) at ($(z\zl x\sp.center)+(\mkDX,-0.38)$);
    \fill[#1, opacity=0.5] (mkpt\zl) circle (2.0pt);
    \ifnum\zl>1
      \pgfmathtruncatemacro{\zp}{\zl-1}%
      \draw[#1,#2,line width=1.1pt, opacity=0.5] (mkpt\zp) -- (mkpt\zl);
    \fi
  }%
}

% Legende: \mkLegende{Nr. (0,1,2,...)}{Farbe}{Linienstil}{Text}
\newcommand{\mkLegende}[4]{%
  \pgfmathsetmacro{\mkLX}{-\mkP/2 + #1*4.6}%
  \pgfmathsetmacro{\mkLY}{-(\mkM+0.8)*\mkH}%
  \draw[#2,#3,line width=1.1pt, opacity=0.5] (\mkLX,\mkLY) -- ++(0.8,0);
  \fill[#2, opacity=0.5] (\mkLX+0.4,\mkLY) circle (2.0pt);
  \node[anchor=west,font=\sffamily\small,text=mkText] at (\mkLX+1.0,\mkLY) {#4};}

% ---------- Abbildung ----------
\begin{figure}[H]
  \captionsetup{font=small,labelfont=bf,labelsep=colon}
  \centering
  \begin{tikzpicture}
    \mkSetup

    % Kopfzeile
    \mkkopfp{0}{Parameter}
    \foreach \c in {1,...,\mkN}{%
      \mkkopf{0}{\c}{Lösung~\c}%
    }

    % Parameterspalte
    \mkparam{1}{Material des Gehäuses}
    \mkparam{2}{Verbindung Gehäuse und Roboter}
    \mkparam{3}{Wärmeabfuhr}
    \mkparam{4}{Fertigung des Gehäuses}
    \mkparam{5}{Energieversorgung/\allowbreak Wandlung}
    \mkparam{6}{Befestigung der Platine im Gehäuse}
    \mkparam{7}{Befestigung der Komponenten}
    \mkparam{8}{Befestigung der Kabel}

    % Material des Gehäuses
    \mkzelle{1}{1}{PLA}
    \mkzelle{1}{2}{PETG}
    \mkzelle{1}{3}{ABS}
    \mkzelle{1}{4}{Aluminium}
    \mkzelle{1}{5}{Acryl}
    \mkzelle{1}{6}{S235JR}

    % Verbindung Gehäuse und Roboter
    \mkzelle{2}{1}{Steckverbindung / Klemmen}
    \mkzelle{2}{2}{Magnetisch}
    \mkzelle{2}{3}{Schraub\-ver\-bin\-dung}
    \mkzelle{2}{4}{Kleben}
    \mkzelle{2}{5}{Einteilig}

    % Wärmeabfuhr
    \mkzelle{3}{1}{Keine}
    \mkzelle{3}{2}{Natürliche Belüftung (Lüftungsschlitze)}
    \mkzelle{3}{3}{Kühlrippen}
    \mkzelle{3}{4}{Lüfter}
    \mkzelle{3}{5}{Wasserkühlung}
    \mkzelle{3}{6}{Flüssiger Stickstoff (LN$_2$)}

    % Fertigung des Gehäuses
    \mkzelle{4}{1}{3D-Druck}
    \mkzelle{4}{2}{CNC-Fräsen}
    \mkzelle{4}{3}{Spritzguss}
    \mkzelle{4}{4}{Manuelles Spanen}

    % Energieversorgung/Wandlung
    \mkzelle{5}{1}{Einzelne Netzteile für 24 V und 5 V}
    \mkzelle{5}{2}{Buck-Converter (24 V $\rightarrow$ 5 V)}
    \mkzelle{5}{3}{Boost-Converter (5 V $\rightarrow$ 24 V)}
    \mkzelle{5}{4}{Trafo mit Gleichrichter (230 V-Anschluss)}
    \mkzelle{5}{5}{Zener-Diode}
    \mkzelle{5}{6}{Spannungsteiler für 5 V}
    \mkzelle{5}{7}{Akku}

    % Befestigung der Platine im Gehäuse
    \mkzelle{6}{1}{Keine}
    \mkzelle{6}{2}{Schrauben ins Material}
    \mkzelle{6}{3}{Schrauben in Einschmelzgewinde}
    \mkzelle{6}{4}{Kleben}
    \mkzelle{6}{5}{Verguss (Epoxidharz)}

    % Befestigung der Komponenten
    \mkzelle{7}{1}{Steckverbindungen (Pins)}
    \mkzelle{7}{2}{Löten (SMD / Through-Hole)}
    \mkzelle{7}{3}{Schraubverbindung}
    \mkzelle{7}{4}{Kleben}

    % Befestigung der Kabel
    \mkzelle{8}{1}{Steckverbindung}
    \mkzelle{8}{2}{Anlöten}
    \mkzelle{8}{3}{Schraubklemmen}
    \mkzelle{8}{4}{Magnetpins}

    % Lösungspfade
    \mkLoesung{mkA}{solid}{-0.25}{1,3,2,1,2,3,1,1}
    \mkLoesung{mkB}{dashed}{0}{2,3,4,1,1,3,2,1}
    \mkLoesung{mkC}{densely dotted}{0.25}{4,3,4,2,2,2,2,2}

    % Legende
    \mkLegende{0}{mkA}{solid}{Lösung 1}
    \mkLegende{1}{mkB}{dashed}{Lösung 2}
    \mkLegende{2}{mkC}{densely dotted}{Lösung 3}
  \end{tikzpicture}
  \caption{Morphologischer Kasten mit drei beispielhaften Lösungskombinationen}
  \label{fig:morphkasten}
\end{figure}

\end{document}